\documentclass{article}
\usepackage{fancyhdr}
\usepackage{graphicx}
\usepackage{makecell}
\usepackage[T1]{fontenc}
\usepackage[margin=0.5in]{geometry}

\title{Titolo dell'esercizio}
\author{DISIM UnivAQ}

% Uso la data della gara 2023
\date{29 Settembre 2023}

\fancyhead[L]{\includegraphics[scale=.125]{header_pic.jpg}}
\fancyhead[R]{\includegraphics[scale=.080]{informatica_univaq.png}}
\fancyhead[C]{\vspace{2pt} \includegraphics[scale=0.9]{logo_nuovo.eps}}

\begin{document}

\maketitle
\thispagestyle{fancy}

\section*{Premessa}

\paragraph*{} Lorem ipsum dolor sit amet, consectetur adipiscing elit. Sed mollis ex sit amet odio tincidunt, vitae feugiat est fermentum. Pellentesque facilisis neque eget sapien tempus, nec pretium mi volutpat. Vestibulum viverra dapibus nibh. Sed pretium rhoncus eros non accumsan. Donec efficitur justo vitae vulputate laoreet. Curabitur efficitur mollis urna. Integer ac nulla ex. Sed quis sem vel massa pretium mattis id eget risus. Sed scelerisque a quam at placerat. Maecenas blandit arcu at convallis vulputate. In convallis metus ac ullamcorper porttitor. In et tortor in nisl vestibulum blandit. Nam leo neque, fringilla quis malesuada sed, iaculis nec mauris. Maecenas porta urna nunc, non congue orci ultrices id. 

\paragraph*{} \textbf{Donec} eget ullamcorper mi. Nulla fringilla hendrerit diam, id feugiat magna pellentesque at. Nullam id viverra urna. Integer vehicula vulputate diam commodo condimentum. Aenean pellentesque, mauris vel eleifend finibus, magna erat fringilla risus, varius venenatis justo eros sodales ligula. Nam ac purus eleifend, luctus quam efficitur, cursus turpis. Mauris in maximus ante. Proin dictum nisi quis auctor efficitur. 



\section*{Problema}

\paragraph*{} Donec eget ullamcorper mi. Nulla fringilla hendrerit diam, id feugiat magna pellentesque at. Nullam id viverra urna. Integer vehicula vulputate diam commodo condimentum. Aenean pellentesque, mauris vel eleifend finibus, magna erat fringilla risus, varius venenatis justo eros sodales ligula.

\begin{enumerate}
    \item un punto;
    \item un altro punto;
    \item un terzo punto.
\end{enumerate}

\textbf{Una tabella, tanto per comodità}.

\begin{center}
    \begin{tabular}{|c|p{10cm}|}
    \hline \rule{0pt}{3ex}
    \makecell{\textit{Prima risposta}} & \\
    \hline \rule{0pt}{3ex}
    \makecell{\textit{Seconda risposta}} &  \\
    \hline \rule{0pt}{3ex}
    \makecell{\textit{Terza risposta}} & \\
    \hline
    \end{tabular}\\
\end{center}


\section*{Suggerimento}

\paragraph*{} Un suggerimento. Questa sezione è \textbf{opzionale}.


\end{document}
